\chapter{Родитель нейтринной линии и квадратичная опора перехода}
\label{ch:v6-spectral-transition-neutrino-line-parent-gate}

\section{Нулевая гипотеза}

Хиггс-разрешённый гейт выделил линию
\begin{equation}
 P_\nu(H)=
 \frac{\widetilde H\widetilde H^\dagger}{H^\dagger H},
 \qquad H\ne0.
 \label{eq:v6-neutrino-normalized-line}
\end{equation}
Теперь проверяется более сильное утверждение: существует ли эта линия как
глобальный непрерывный и $SU(2)$-ковариантный объект родителя, включая
точки $H=0$? Если нет, требуется определить минимальный объект, который
остаётся регулярным без введения правого нейтрино.

Это различие существенно для спектральной парадигмы. Глобальная линия
была бы постоянным внутренним «волокном». Регулярный оператор, ранг
которого меняется вместе с фазой, является квантом перехода, а не
геометрической верёвкой.

\section{Глобальный проектор запрещён симметрией}

Коммутант фундаментального действия $SU(2)$ на $L_L\simeq\mathbb C^2$
состоит только из скаляров:
\begin{equation}
 \operatorname{End}_{SU(2)}(\mathbb C^2)=\mathbb C I_2.
\end{equation}
Поэтому invariant-проекторы имеют только ранги ноль и два.

Точка $H=0$ фиксирована всей группой $SU(2)$. Если бы
$P_\nu(H)$ имел непрерывное эквивариантное продолжение в эту точку, его
значение $P_\nu(0)$ должно было бы быть инвариантным проектором ранга
один. Такого проектора нет.

Прямой предел показывает ту же проблему. Для путей
\begin{equation}
 H_1(r)=(r,0),\qquad H_2(r)=(0,r)
\end{equation}
проекторы не зависят от $r>0$, но различны, а их расстояние Фробениуса
равно $\sqrt2$. Значит, единого предела при $H\to0$ нет.

\begin{theorem}[No-go глобальной нейтринной линии]
Не существует непрерывного $SU(2)$-эквивариантного продолжения
нормированного проектора \eqref{eq:v6-neutrino-normalized-line} через
$H=0$. Следовательно, нейтринная линия не является глобальным
проектирующим слагаемым замороженного родителя $H_{15}$.
\end{theorem}

Этот вывод согласуется с конечной диаграммной границей: вершины задают
фермионные мультиплеты, а рёбра --- юкавские связи; на $H_{15}$ нет
вершины $\nu_R$ и степени-один нейтринного ребра
\cite{neutrino-parent-krajewski1998,
neutrino-parent-chamseddine-connes-marcolli2007}.

\section{Регулярный квадратичный носитель}

Нормировка не нужна для задания самой опоры перехода. Определим
\begin{equation}
 W_\nu(H)=\widetilde H\widetilde H^\dagger,
 \qquad
 B_\nu(H)=\widetilde H\widetilde H^T.
 \label{eq:v6-neutrino-unnormalized-parent}
\end{equation}
Первый объект является эрмитовым носителем, второй --- симметричным
паринговым тензором. Они полиномиальны второй степени, ковариантны и
удовлетворяют
\begin{equation}
 \rank W_\nu(H)=1\quad(H\ne0),
 \qquad W_\nu(0)=0.
\end{equation}
Нормированный проектор восстанавливается только на ненулевой области:
\begin{equation}
 P_\nu(H)=\frac{W_\nu(H)}{\Tr W_\nu(H)}.
\end{equation}

Машинный аудит дал остатки ковариантности не выше
$6.23\cdot10^{-15}$ и точную квадратичную однородность с остатком
$2.89\cdot10^{-17}$. В отличие от проектора, $W_\nu$ имеет единственное
продолжение в ноль: нулевой оператор. Цена регулярности состоит в смене
ранга $1\to0$.

\section{Связь с оператором Вайнберга}

На $H_{15}$ нет правого нейтрино, поэтому степень-один дираковская масса
невозможна. Первая стандартная gauge-инвариантная возможность использует
две лептонные и две хиггсовские ноги, то есть оператор размерности пять
\cite{neutrino-parent-weinberg1979}. Его внутреннее направление
согласуется с $B_\nu(H)$, но из существования квадратичного тензора ещё
не следуют:
\begin{enumerate}
 \item коэффициент оператора и подавляющий масштаб;
 \item семейная симметричная матрица;
 \item механизм нарушения лептонного числа;
 \item совместный нелинейный функционал, рождающий ненулевой $H$.
\end{enumerate}

Семейный голономный проектор $P_0$ может условно дать
\begin{equation}
 P_0\otimes W_\nu(H),
\end{equation}
который имеет ранг один при $H\ne0$ и ранг ноль при $H=0$. Голономия
фиксирует семейное направление, но не коэффициент паринга.

\begin{proposition}[Правильный спектральный примитив нейтринной ветви]
Минимальным глобально регулярным объектом является не нормированная
нейтринная линия, а ненормированный хиггс-квадратичный ковариант
$W_\nu(H)$ либо его паринговая версия $B_\nu(H)$. Он описывает рождение
ранг-один опоры из нулевой опоры при переходе в фазу $H\ne0$.
\end{proposition}

\begin{nogocriterion}[Запрет деления на параметр порядка]
Нельзя считать $P_\nu(H)$ фундаментальным глобальным полем, если
родитель допускает $H=0$. Деление на $H^\dagger H$ удаляет физически
значимую смену ранга и создаёт ложную сингулярную «вечную линию».
\end{nogocriterion}

\section{Вердикт и следующий гейт}

Нормированная линия закрыта отрицательно как глобальный parent-объект.
Квадратичная опора закрыта положительно как канонический допустимый
ковариант поля Хиггса. Но массовая динамика остаётся открытой.

Следующий гейт
\texttt{version6\_spectral\_transition\_weinberg\_pairing\_parent\_gate}
должен проверить, выводит ли существующая спектральная архитектура сам
паринговый оператор, его семейный тензор и коэффициент, либо они остаются
новыми внешними данными.

\begin{protocol}
\begin{enumerate}
 \item постоянный gauge-инвариантный проектор ранга один:
 \textbf{невозможен};
 \item непрерывное продолжение $P_\nu$ через $H=0$:
 \textbf{невозможно};
 \item регулярный квадратичный носитель $W_\nu$:
 \textbf{существует};
 \item ранг $W_\nu$: \textbf{$0$ при $H=0$, $1$ при $H\ne0$};
 \item степень-один нейтринное ребро на $H_{15}$:
 \textbf{отсутствует};
 \item оператор Вайнберга: \textbf{структурно допустим};
 \item его коэффициент, семейный тензор и масштаб:
 \textbf{не выведены};
 \item физическое замыкание: \textbf{не достигнуто}.
\end{enumerate}
\end{protocol}

Машинный сертификат сохранён в
\path{s2t_v6_spectral_transition_neutrino_line_parent_gate.py} и
\path{s2t_v6_spectral_transition_neutrino_line_parent_gate_results.json}.

\begin{thebibliography}{9}
\bibitem{neutrino-parent-krajewski1998}
T.~Krajewski,
\emph{Classification of Finite Spectral Triples},
Journal of Geometry and Physics 28 (1998), 1--30;
arXiv:hep-th/9701081.

\bibitem{neutrino-parent-chamseddine-connes-marcolli2007}
A.~H.~Chamseddine, A.~Connes, M.~Marcolli,
\emph{Gravity and the Standard Model with Neutrino Mixing},
Advances in Theoretical and Mathematical Physics 11 (2007), 991--1089;
arXiv:hep-th/0610241.

\bibitem{neutrino-parent-weinberg1979}
S.~Weinberg,
\emph{Baryon- and Lepton-Nonconserving Processes},
Physical Review Letters 43 (1979), 1566--1570.
\end{thebibliography}
